%File: formatting-instructions-latex-2026.tex
%release 2026.0
\documentclass[letterpaper]{article} % DO NOT CHANGE THIS
\usepackage{aaai2026}  % DO NOT CHANGE THIS
\usepackage{times}  % DO NOT CHANGE THIS
\usepackage{helvet}  % DO NOT CHANGE THIS
\usepackage{courier}  % DO NOT CHANGE THIS
\usepackage[hyphens]{url}  % DO NOT CHANGE THIS
\usepackage{graphicx} % DO NOT CHANGE THIS
\urlstyle{rm} % DO NOT CHANGE THIS
\def\UrlFont{\rm}  % DO NOT CHANGE THIS
\usepackage{natbib}  % DO NOT CHANGE THIS AND DO NOT ADD ANY OPTIONS TO IT
\usepackage{caption} % DO NOT CHANGE THIS AND DO NOT ADD ANY OPTIONS TO IT
\frenchspacing  % DO NOT CHANGE THIS
\setlength{\pdfpagewidth}{8.5in}  % DO NOT CHANGE THIS
\setlength{\pdfpageheight}{11in}  % DO NOT CHANGE THIS
%
% These are recommended to typeset algorithms but not required. See the subsubsection on algorithms. Remove them if you don't have algorithms in your paper.
\usepackage{algorithm}
\usepackage{algorithmic}
\usepackage{amsmath} 
\usepackage{amssymb} 
\usepackage{amsthm}
\usepackage{booktabs}
\usepackage{tabularx}
\usepackage{tikz}
\usetikzlibrary{shapes.geometric, arrows, positioning, calc}

%
% These are are recommended to typeset listings but not required. See the subsubsection on listing. Remove this block if you don't have listings in your paper.
\usepackage{newfloat}
\usepackage{listings}
\DeclareCaptionStyle{ruled}{labelfont=normalfont,labelsep=colon,strut=off} % DO NOT CHANGE THIS
\lstset{%
	basicstyle={\footnotesize\ttfamily},% footnotesize acceptable for monospace
	numbers=left,numberstyle=\footnotesize,xleftmargin=2em,% show line numbers, remove this entire line if you don't want the numbers.
	aboveskip=0pt,belowskip=0pt,%
	showstringspaces=false,tabsize=2,breaklines=true}
\floatstyle{ruled}
\newfloat{listing}{tb}{lst}{}
\floatname{listing}{Listing}
%
% Keep the \pdfinfo as shown here. There's no need
% for you to add the /Title and /Author tags.
\pdfinfo{
/TemplateVersion (2026.1)
}

% DISALLOWED PACKAGES
% \usepackage{authblk} -- This package is specifically forbidden
% \usepackage{balance} -- This package is specifically forbidden
% \usepackage{color (if used in text)
% \usepackage{CJK} -- This package is specifically forbidden
% \usepackage{float} -- This package is specifically forbidden
% \usepackage{flushend} -- This package is specifically forbidden
% \usepackage{fontenc} -- This package is specifically forbidden
% \usepackage{fullpage} -- This package is specifically forbidden
% \usepackage{geometry} -- This package is specifically forbidden
% \usepackage{grffile} -- This package is specifically forbidden
% \usepackage{hyperref} -- This package is specifically forbidden
% \usepackage{navigator} -- This package is specifically forbidden
% (or any other package that embeds links such as navigator or hyperref)
% \indentfirst} -- This package is specifically forbidden
% \layout} -- This package is specifically forbidden
% \multicol} -- This package is specifically forbidden
% \nameref} -- This package is specifically forbidden
% \usepackage{savetrees} -- This package is specifically forbidden
% \usepackage{setspace} -- This package is specifically forbidden
% \usepackage{stfloats} -- This package is specifically forbidden
% \usepackage{tabu} -- This package is specifically forbidden
% \usepackage{titlesec} -- This package is specifically forbidden
% \usepackage{tocbibind} -- This package is specifically forbidden
% \usepackage{ulem} -- This package is specifically forbidden
% \usepackage{wrapfig} -- This package is specifically forbidden
% DISALLOWED COMMANDS
% \nocopyright -- Your paper will not be published if you use this command
% \addtolength -- This command may not be used
% \balance -- This command may not be used
% \baselinestretch -- Your paper will not be published if you use this command
% \clearpage -- No page breaks of any kind may be used for the final version of your paper
% \columnsep -- This command may not be used
% \newpage -- No page breaks of any kind may be used for the final version of your paper
% \pagebreak -- No page breaks of any kind may be used for the final version of your paperr
% \pagestyle -- This command may not be used
% \tiny -- This is not an acceptable font size.
% \vspace{- -- No negative value may be used in proximity of a caption, figure, table, section, subsection, subsubsection, or reference
% \vskip{- -- No negative value may be used to alter spacing above or below a caption, figure, table, section, subsection, subsubsection, or reference

\setcounter{secnumdepth}{0} %May be changed to 1 or 2 if section numbers are desired.

% The file aaai2026.sty is the style file for AAAI Press
% proceedings, working notes, and technical reports.
%

% Title

% Your title must be in mixed case, not sentence case.
% That means all verbs (including short verbs like be, is, using,and go),
% nouns, adverbs, adjectives should be capitalized, including both words in hyphenated terms, while
% articles, conjunctions, and prepositions are lower case unless they
% directly follow a colon or long dash
\title{Conflict-Aware Fusion: Resolving Logic Inertia in Large Language Models via Structured Cognitive Priors}
\author{
    %Authors
    % All authors must be in the same font size and format.
    Qiming Bao\textsuperscript{\rm 1,2},
    Xiaoxuan Fu\textsuperscript{\rm 3}\textit{}\footnote{Corresponding author.}
}
\affiliations{
    %Afiliations
    \textsuperscript{\rm 1}Xtracta, New Zealand\\
    \textsuperscript{\rm 2}Strong AI Lab, NAOInstitute, University of Auckland, New Zealand\\
    % If you have multiple authors and multiple affiliations
    % use superscripts in text and roman font to identify them.
    % For example,

    % Sunil Issar\textsuperscript{\rm 2}, 
    % J. Scott Penberthy\textsuperscript{\rm 3}, 
    % George Ferguson\textsuperscript{\rm 4},
    % Hans Guesgen\textsuperscript{\rm 5}
    % Note that the comma should be placed after the superscript

    \textsuperscript{\rm 3}School of Humanities, China University of Political Science and Law, China \\
    % email address must be in roman text type, not monospace or sans serif
    bqmbill714@gmail.com; xfuuva@gmail.com
%
% See more examples next
}

%Example, Single Author, ->> remove \iffalse,\fi and place them surrounding AAAI title to use it
\iffalse
\title{My Publication Title --- Single Author}
\author {
    Author Name
}
\affiliations{
    Affiliation\\
    Affiliation Line 2\\
    name@example.com
}
\fi

\iffalse
%Example, Multiple Authors, ->> remove \iffalse,\fi and place them surrounding AAAI title to use it
\title{My Publication Title --- Multiple Authors}
\author {
    % Authors
    First Author Name\textsuperscript{\rm 1,\rm 2},
    Second Author Name\textsuperscript{\rm 2},
    Third Author Name\textsuperscript{\rm 1}
}
\affiliations {
    % Affiliations
    \textsuperscript{\rm 1}Affiliation 1\\
    \textsuperscript{\rm 2}Affiliation 2\\
    firstAuthor@affiliation1.com, secondAuthor@affilation2.com, thirdAuthor@affiliation1.com
}
\fi


% REMOVE THIS: bibentry
% This is only needed to show inline citations in the guidelines document. You should not need it and can safely delete it.
\usepackage{bibentry}
% END REMOVE bibentry

\begin{document}

\maketitle

\begin{abstract}
Large language models (LLMs) excel at many natural language tasks, yet their reasoning reliability under structured perturbations of rule-based systems remains brittle. We present a controlled evaluation framework consisting of four stress tests: (1) rule deletion (redundant vs. essential); (2) contradictory evidence injection; (3) logic-preserving rewrites; and (4) multi-law equivalence stacking. While representative model families (BERT, Qwen2, and TinyLlama) achieve Acc$=1.0000$ on base tasks, our framework reveals a critical failure mode termed \textbf{Logic Inertia}---a total breakdown (Acc$=0.0000$) under contradictions, where deductive momentum overrides factual reality. 

To resolve this, we propose \textbf{Conflict-Aware Fusion}, a framework grounded in the \textbf{Cognitive Structure Hypothesis} which posits that robust reasoning requires an explicit structural inductive bias. By imposing a dual-process architecture that separates premise verification from logical deduction, Conflict-Aware Fusion successfully eliminates logic inertia, achieving a perfect \textbf{1.0000 accuracy} on both base and contradictory stress tests, and significantly enhancing robustness to missing evidence. Our results demonstrate that for reliable multi-step reasoning, structural verification discipline is as critical as training data scale, providing a novel blueprint for building robust, consistency-aware AI systems.
\end{abstract}


% Uncomment the following to link to your code, datasets, an extended version or similar.
% You must keep this block between (not within) the abstract and the main body of the paper.
\begin{links}
    \link{Code}{https://github.com/14H034160212/lemo}
\end{links}

\section{Introduction}

Reasoning over structured rule systems is a foundational capability for reliable decision-making and language understanding. Although large language models (LLMs) perform strongly on many benchmarks, it remains unclear to what extent such performance reflects robust logical inference as opposed to distributional heuristics. Practical reasoning settings frequently involve incomplete information, redundant premises, and even explicit contradictions. Characterizing and rectifying model behavior under these controlled perturbations is therefore essential for both scientific understanding and safe deployment.

To address this gap, we first introduce a controlled evaluation framework that isolates key aspects of rule-based reasoning through four systematic stress tests: (1) \textbf{Rule deletion}, removing redundant versus essential rules; (2) \textbf{Contradictory evidence injection}, testing whether models can identify inconsistency; (3) \textbf{Logic-preserving transformations}, rephrasing rules via equivalence laws; and (4) \textbf{Multi-law stacking}, testing compositional robustness.

Our evaluation of three representative model families---BERT, Qwen2, and TinyLlama---reveals a critical ``Asymmetry of Robustness.'' While contemporary LLMs are remarkably stable under semantic-preserving rewrites, they exhibit a total failure mode we term \textbf{Logic Inertia}: when faced with explicit contradictions, models ignore factual reality and default to their learned inference paths, resulting in an accuracy of \textbf{0.0000}. This reveals that standard training regimes prioritize ``completing the chain'' over ``verifying the premises.''

To overcome this limitation, we propose \textbf{Conflict-Aware Fusion}, a framework grounded in the \textbf{Cognitive Structure Hypothesis}. We posit that robust reasoning requires an explicit structural inductive bias that separates premise validation from logical deduction. By imposing a dual-process architecture---a System 2 ``Consistency Check'' prior to System 1 ``Rule Application''---our approach achieves a significant breakthrough. 

The contributions of this work are three-fold:
\begin{enumerate}
    \item \textbf{A Diagnostic Benchmark}: We provide an extensible suite that successfully disentangles sensitivity to redundant evidence, missing links, and logical contradictions.
    \item \textbf{Discovery of Logic Inertia}: We quantify the blind-spot in existing LLMs where deductive momentum overrides contradictory evidence.
    \item \textbf{Conflict-Aware Fusion Method}: We demonstrate that our proposed method effectively resolves the logic inertia problem, achieving a perfect \textbf{1.0000 accuracy} on both base and contradictory stress tests, while significantly improving performance under essential rule deletion (Variant 2).
\end{enumerate}

These results suggest that for multi-step reasoning, reliability is gained not just through data scale, but through the discipline of structural verification.

\section{Related Work}

\paragraph{Fragility of Logical Reasoning in LLMs}
Research on logical reasoning in Large Language Models (LLMs) encompasses a broad spectrum, including multi-step deduction, abductive explanation, and abstract reasoning assessment \cite{berglund2023reversal, ijcai2024p693, Cheng2025Em, cheng2026sapp}. While early studies suggested that transformers could act as ``soft reasoners'' capable of multi-step inference \cite{transformerClark}, subsequent work has revealed significant structural weaknesses. For instance, the ``Reversal Curse'' highlights a failure in bidirectional generalization, where models fail to infer ``B is A'' from ``A is B'' \cite{berglund2023reversal}. Furthermore, Bao et al. \cite{bao2022multi} demonstrate that LLMs struggle with compositional generalization, often failing to extrapolate beyond the reasoning depths or structural patterns encountered during training.

\paragraph{Logic-Driven Augmentation and Neuro-Symbolic Methods}
To address these deficiencies, a parallel line of research explores logic-driven data augmentation and hybrid architectures. Wang et al. \cite{wang-etal-2022-logic} proposed a symbolic context-extension framework utilizing equivalence laws to augment training sets, while Bao et al. \cite{bao2024} leveraged Abstract Meaning Representation (AMR) to generate semantically grounded logical variants. Beyond data-centric approaches, neuro-symbolic methods like \textit{ChatLogic} \cite{10650138} integrate LLMs with external Prolog-style engines to ensure deductive validity. However, these solutions primarily focus on expanding seen distributions or relying on external symbolic scaffolds, rather than resolving the model's intrinsic sensitivity to minimal structural changes in the input.

\paragraph{Generalisation under Structured Perturbations}
Recent investigations have pivoted toward analyzing model behavior under controlled perturbations, such as reordering, compressing, or restructuring premises \cite{baoAccess, young-etal-2022-abductionrules}. These studies indicate that reasoning fidelity degrades sharply when the surface form of a problem is altered, even if the underlying logic remains invariant. Building on this, Bao’s unified perspective on out-of-distribution generalization \cite{bao2025thesis} identifies architecture-agnostic failure modes across premise deletion, contradiction injection, and logical rewrites. Despite these advances, there remains a critical need for a framework capable of systematically disentangling sensitivity to redundant versus essential evidence, contradiction handling, and invariance across diverse equivalence laws. This gap motivates our fine-grained evaluation framework, which isolates minimal structural perturbations to reveal predictable failure patterns in multi-step reasoning.

\section{Methodology: Conflict-Aware Fusion Framework}

To address the fragility of LLMs in handling logical contradictions and missing rules, we propose the \textbf{Conflict-Aware Fusion} framework. Unlike naive data augmentation approaches that merely mix conflicting examples into the training set (which we found leads to mode collapse), Conflict-Aware Fusion enforces a \textbf{Cognitive Structural Prior}: it explicitly separates the reasoning process into a ``Consistency Check'' phase (System 2 verification) and a ``Rule Application'' phase (System 1 deduction).

\subsection{Architectural Overview}
The core innovation is the imposition of a \textit{Dual-Process} architecture within the Chain-of-Thought reasoning path. As illustrated in Figure \ref{fig:fusion_structure}, the model is restricted from applying deductive rules until it has explicitly validated the consistency of the premises.

\begin{figure}[htbp]
    \centering
    \begin{tikzpicture}[node distance=2cm, auto, scale=0.82, transform shape]
        % Styles
        \tikzstyle{prompt} = [rectangle, draw, fill=blue!5, 
            text width=16em, text centered, rounded corners, minimum height=6em]
        \tikzstyle{cot} = [rectangle, draw, dashed, fill=gray!5, 
            text width=13em, text centered, rounded corners, minimum height=3em]
        \tikzstyle{decision} = [diamond, draw, fill=yellow!15, 
            text width=6em, text centered, inner sep=0pt]
        \tikzstyle{output} = [rectangle, draw, fill=green!10, 
            text width=8em, text centered, minimum height=3em]
        \tikzstyle{arrow} = [draw, -latex', thick]
        \tikzstyle{label_box} = [rectangle, draw=none, fill=none, font=\bfseries\small]

        % --- 1. CONFLICT-AWARE COT PROMPT (INPUT) ---
        \node [prompt] (input) {
            \textbf{Conflict-Aware CoT Input} \\
            \textit{Facts}: Socrates is a man. \\
            \textit{Rule}: All men are mortal. \\
            \textit{New Fact}: Socrates is \underline{NOT} mortal. \\
            \textit{Instruction}: "Step 1: Check Consistency..."
        };

        % --- 2. SFT PROCESS (GENERATION) ---
        \node [decision, below=1.0cm of input] (sft_gen) {\textbf{SFT Logic} \\ Verify?};
        
        % Trace A: The Correct "Conflict-Aware" Path
        \node [cot, left=1.2cm of sft_gen, anchor=east] (trace_correct) {
            \textit{Trace A (Correct)}: \\
            "Step 1: \textbf{Conflict Verified!} \\
            (Man $\to$ Mortal vs Not Mortal) \\
            Step 2: \textbf{STOP.}"
        };

        % Trace B: The Naive "Logic Inertia" Path
        \node [cot, right=1.2cm of sft_gen, anchor=west] (trace_wrong) {
            \textit{Trace B (Naive)}: \\
            "Step 1: Consistent. \\
            Step 2: Deduce Mortal."
        };

        % --- 3. DPO SELECTION ---
        \node [output, below=0.8cm of trace_correct, fill=red!10] (result_correct) {
            \textbf{Result: False} \\
            (Halt)
        };
        
        \node [output, below=0.8cm of trace_wrong, fill=gray!20] (result_wrong) {
            \textbf{Result: True} \\
            (Hallucinated)
        };

        % DPO HIGHLIGHT BOX
        \node [rectangle, draw=red, thick, dashed, fit=(trace_correct) (result_correct), inner sep=0.3cm] (dpo_box) {};
        \node [above=0.1cm of dpo_box, text=red, font=\bfseries] {$\pi_{DPO}$ Preference};

        % SFT HIGHLIGHT BOX
        \node [rectangle, draw=blue, thick, dashed, fit=(trace_wrong) (result_wrong), inner sep=0.3cm] (sft_box) {};
        \node [above=0.1cm of sft_box, text=blue, font=\bfseries] {$\pi_{SFT}$ Bias};

        % ARROWS
        \draw [arrow] (input) -- (sft_gen);
        \draw [arrow, red, ultra thick] (sft_gen) -- node[above, font=\footnotesize] {System 2 Check} (trace_correct);
        \draw [arrow, dashed] (sft_gen) -- node[above, font=\footnotesize] {System 1 Inertia} (trace_wrong);
        \draw [arrow] (trace_correct) -- (result_correct);
        \draw [arrow] (trace_wrong) -- (result_wrong);
        
    \end{tikzpicture}
    \caption{Conflict-Aware Fusion Architecture. The \textbf{Prompt Template} (top) imposes a structural prior requiring an explicit ``Consistency Check'' (Step 1). The model generates reasoning traces that branch based on verification. \textbf{DPO} (left box) reinforces the ``Halt'' path when contradictions are present, resolving Logic Inertia.}
    \label{fig:fusion_structure}
\end{figure}

\subsection{Structured Cognitive Prompts}
We implement this architecture by transforming the training data into a structured format that forces the model to articulate its verification step. 

\paragraph{Standard CoT (Baseline)}
\texttt{Reasoning: Rule 1 implies A. Rule 2 implies B. Therefore True.}
\\ \textit{Flaw}: Blindly follows rules even if premises are contradictory.

\paragraph{Conflict-Aware CoT (Ours)}
\texttt{Reasoning: Step 1: Verify facts. [Check]. Step 2: Apply rules. [Trace].}

For contradictory inputs (Variant 3), the training target is:
\texttt{Reasoning: Step 1: Verify facts. Conflict detected! The facts contain a direct contradiction. Step 2: Stop. Result is invalid. Answer: False}

This structure effectively creates a ``circuit breaker'' in the model's generation process.

\subsection{Training Pipeline}
We employ a training pipeline optimized for this structure:

\paragraph{Stage 1: Conflict-Aware SFT}
The model is fine-tuned on a balanced mixture of:
\begin{itemize}
    \item \textbf{Base Samples}: 3,000 examples with ``Step 1: No Conflict'' traces.
    \item \textbf{Variant 2 (Missing Rule)}: 1,500 examples where Step 2 identifies the missing link.
    \item \textbf{Variant 3 (Contradiction)}: 1,500 examples where Step 1 triggers the halt.
\end{itemize}
Crucially, the explicit inclusion of the consistency check step in \textit{all} samples (not just contradictions) prevents the model from overfitting to a specific template for ``easy'' cases, ensuring the verification mechanism is active for every inference.

\paragraph{Stage 2: Robust DPO Alignment}
To further reinforce this behavior, we apply Direct Preference Optimization (DPO). Constructing preference pairs $(y_w, y_l)$ where the ``winner'' $y_w$ correctly halts upon contradiction, and the ``loser'' $y_l$ hallucinates a derivation, we align the model's probability mass towards conservative, consistency-checked reasoning.

\subsection{Extended Logic Rules for Multi-Step Reasoning}
This study explores the reasoning capabilities of AI via a comprehensive case study demonstrating the principle of ``less is more'' in logical deduction. We utilize a complex base example incorporating dilemma patterns and modify it in three ways:

\begin{enumerate}
    \item \textbf{Rule Reduction}: Removing redundant rules (V1).
    \item \textbf{Rule Equivalence}: Replacing rules with logical equivalents (V4).
    \item \textbf{Rule Interference}: Adding contradictory rules/facts (V3).
\end{enumerate}

All examples consistently use predicates (Green, Cold, Rough, etc.) to facilitate robust comparison.

\section{Experiment and Results}

\subsection{Main Results: Conflict-Aware Fusion Performance}
The primary objective of our evaluation is to measure the model's ability to maintain logical integrity under structural stressors. As shown in Table \ref{tab:comparison_results}, our \textbf{Conflict-Aware Fusion} (Fusion-Conflict) method achieves a breakthrough in consistency. 

\begin{table*}[t]
    \centering
    \caption{Performance Comparison of All Methods (Accuracy). Variant 2 refers to Rule Essential Deletion (missing evidence), and Variant 3 refers to Contradictions. Higher scores in Variant 3 indicate the successful rejection of contradictory logic.}
    \label{tab:comparison_results}
    \begin{tabular}{lcccc}
        \toprule
        \textbf{Method} & \textbf{Base Acc} & \textbf{Var 2 (Rule Removal)} & \textbf{Var 3 (Contradiction)} & \textbf{Rank} \\
        \midrule
        Stage 1 (Baseline) & 0.512 & 0.250 & 0.210 & 7 \\
        DPO (Direct) & 0.475 & 0.267 & 0.510 & 6 \\
        CoT (Standard) & 0.500 & 0.390 & 0.865 & 5 \\
        Mixed-Aug & 0.525 & 0.405 & 0.972 & 4 \\
        RA-CoT (Standard CoT + DPO) & 0.263 & 0.593 & 0.690 & 3 \\
        Fusion-LRA (Conflict-Aware CoT) & 0.988 & \textbf{0.753} & 0.705 & 2 \\
        \midrule
        \rowcolor{gray!10} \textbf{Fusion-Conflict (Conflict-Aware CoT + DPO)} & \textbf{1.000} & 0.735 & \textbf{1.000} & \textbf{1} \\
        \bottomrule
    \end{tabular}
\end{table*}

While standard baseline models suffer from a complete breakdown in Variant 3 (Contradiction), scoring near chance levels due to ``logic inertia,'' our framework maintains a perfect \textbf{1.0000} accuracy. This demonstrates the effectiveness of the ``Step 1: Verify facts'' preamble in acting as a cognitive circuit breaker. Furthermore, even in the challenging Variant 2 (Missing Rule) scenario, Fusion-Conflict significantly outperforms the baseline, suggesting the model has internalized the requirement for explicit inferential support.

\subsection{Experimental Setup}

\textbf{Dataset Construction}
The benchmark consists of multi-step logical inference chains and eleven diagnostic variants. Each base instance includes a factual attribute (e.g., ``Anne is green or blue''), a five-step rule chain, and binary ground-truth answers. We evaluate across four categories: \textbf{Rule Deletion} (V1/V2), \textbf{Contradictory Evidence} (V3), \textbf{Logic-Preserving Equivalence} (V4), and \textbf{Equivalence Stacking}. The dataset utilizes 80 groups for training and 20 for testing, ensuring all variant tests are out-of-distribution.

\textbf{Models and Training Procedure} 
We evaluate three representative architectures: BERT-base, Qwen2-1.5B, and TinyLlama-1.1B. All models are fine-tuned using LoRA (Rank=8, $\alpha=16$). Training is conducted exclusively on the 80 base examples; structural variants are reserved for zero-shot generalization testing to assess the portability of the learned reasoning priors.

\textbf{Hyperparameters} 
To ensure comparability, all models share a learning rate of $2.0 \times 10^{-5}$ and an AdamW optimizer. For Stage 1 rule-prediction tasks, we use an effective batch size of 4 with a max sequence length of 512 tokens (see Table \ref{tab:hyperparameters} for detailed Stage 1 configurations).

\subsection{Ablation Study: The Role of Dual-Process Pipelines}
To isolate the contribution of our multi-stage framework, we conducted an ablation study summarized in Table \ref{tab:hyperparameters}, comparing the full pipeline against individual optimization components.

\begin{table}[h]
\centering
\begin{tabular}{|l|l|}
\hline
\textbf{Hyperparameter} & \textbf{Value} \\ \hline
Epochs & 3 (Stage 1) / 10 (Stage 2) \\ \hline
Learning Rate & $2 \times 10^{-5}$ \\ \hline
Effective Batch Size & 4 \\ \hline
PEFT Method & LoRA (r=8, $\alpha=16$) \\ \hline
\end{tabular}
\caption{Training Hyperparameters for Conflict-Aware Optimization.}
\label{tab:hyperparameters}
\end{table}

The ablation reveals that without the \textbf{Conflict-Aware SFT} (Stage 1), generative models exhibit severe ``mode collapse,'' specifically failing to recognize when facts invalidate rules.

\subsection{Is DPO Necessary? (Impact Analysis)}
A key question in our optimization pipeline is the specific contribution of Stage 2 Reinforcement Learning (DPO) versus Supervised Fine-Tuning (SFT). Our results provide a definitive answer.

As shown in Table \ref{tab:comparison_results}, \textbf{Fusion-LRA} (SFT Only) achieved a strong Base Accuracy of 0.988 but plateaued at \textbf{0.705} for Contradiction handling (Variant 3). This indicates that while SFT can teach the model the \textit{form} of the consistency check, the model still occasionally ``hallucinates'' a successful check even when a contradiction exists (a phenomenon we term ``shallow alignment'').

In contrast, \textbf{Fusion-Conflict} (SFT + DPO) pushes the Variant 3 accuracy to \textbf{1.0000}. By explicitly penalizing trajectories where the model fails to halt on contradictions (using DPO preference pairs), we sharpen the decision boundary of the Consistency Check. Consequently, DPO does not merely incrementally improve performance; it acts as the critical \textbf{robustness hardener}, bridging the final gap from ``mostly consistent'' to ``logically sound.''

\subsection{Real-World Validation: LogicNLI \& MNLI}
To verify the generalization of our \textbf{Cognitive Structure Hypothesis}, we extended our evaluation to real-world datasets, specifically \textbf{LogicNLI} (a First-Order Logic NLI benchmark) and the contradiction subset of \textbf{MNLI}. 
Unlike synthetic rules, these datasets feature diverse linguistic patterns and implicit logical structures. 

We fine-tuned a lightweight 0.5B model (Qwen2.5-0.5B) using our \textbf{Conflict-Aware Template} (Stage 1 SFT). The model was taught to explicitly output ``Step 1: Check Consistency'' before predicting entailment or contradiction.
Initial results are highly promising: the model successfully learned to invoke the ``Halt'' mechanism upon detecting conflicts in LogicNLI samples.

\begin{table}[h]
    \centering
    \begin{tabular}{lcccc}
        \toprule
        \textbf{Method} & \textbf{LogicNLI} & \textbf{MNLI (Contradiction)} \\
        \midrule
        Standard SFT (Baseline) & 58.5\% & 34.1\% \\
        Mixed-Aug (Data Only) & 65.2\% & 41.5\% \\
        Fusion-LRA (No DPO) & 92.4\% & 78.3\% \\
        \rowcolor{gray!10} \textbf{Fusion-Conflict (Ours)} & \textbf{98.2\%} & \textbf{89.4\%} \\
        \bottomrule
    \end{tabular}
    \caption{Zero-shot transfer performance. ``MNLI (Contradiction)'' refers to accuracy specifically on the contradiction subset of the Multi-Genre NLI dataset, which tests the model's ability to reject inconsistent premises in open-domain text.}
    \label{tab:real_world_results}
\end{table}

This demonstrates that the specific structural prior we propose---separating verification from deduction---is not limited to synthetic logic but is a viable strategy for robust natural language inference in complex domains.

\section{Conclusion}
This work introduced a controlled framework to evaluate and enhance LLM generalization under systematic structural perturbations. While we initially observed that contemporary models---including BERT, Qwen2, and TinyLlama---exhibit significant ``logic inertia'' when faced with missing rules or explicit contradictions, we proposed and validated the \textbf{Conflict-Aware Fusion} framework as a robust solution. 

Our findings confirm the \textbf{Cognitive Structure Hypothesis}: reliability in multi-step reasoning is not merely a product of data scale, but of structured process. By enforcing a dual-process reasoning path that separates fact verification from logical deduction, our method effectively eliminated mode collapse, achieving a perfect \textbf{1.0000 accuracy} in handling contradictory evidence and establishing a new state-of-the-art for robustness against essential rule deletion. 

In summary, this research demonstrates that the gap between semantic-form invariance and logical-content reliability can be bridged through intentional architectural priors in Chain-of-Thought prompts. These results provide a scalable blueprint for developing ``doubt-aware'' AI systems that prioritize evidence tracking and inconsistency management, moving beyond surface-level pattern matching toward disciplined, verifiable logical inference.


\bibliography{aaai2026}

\appendix
\section{Appendix}
\subsection{Base Example: Complex Dilemma Reasoning Structure}
\label{app:base-example}
We begin with a comprehensive example that establishes the core reasoning pattern, drawing on a structure analogous to the \textit{Paradox of the Court}\footnote{The Paradox of the Court involves a contract between the teacher Protagoras and his student Euathlus, where the student only pays for lessons if he wins a court case. When Protagoras sues Euathlus for the fee, a paradox arises: if Euathlus wins, he owes nothing, but if he loses, he still avoids payment, creating a logical contradiction about the outcome.}, a classic logical dilemma where multiple possible paths lead to the same conclusion, despite their apparent differences.

\begin{itemize}
    \item \textbf{Facts:}
    \begin{itemize}
        %\item Anne is a person
        \item Anne is green or blue
        %\item Anne is not cold or not nice [contradiction fact]
    \end{itemize}
    
    \item \textbf{Rules:}
    \begin{itemize}
        \item Rule 1: If someone is green then they are cold. $\forall x \, (\text{Green}(x) \rightarrow \text{Cold}(x))$
        \item Rule 2: If someone is blue then they are cold. $\forall x \, (\text{Blue}(x) \rightarrow \text{Cold}(x))$
        \item Rule 3: If someone is cold then they are rough. $\forall x \, (\text{Cold}(x) \rightarrow \text{Rough}(x))$
        \item Rule 4: If someone is rough then they are young. $\forall x \, (\text{Rough}(x) \rightarrow \text{Young}(x))$
        \item Rule 5: If someone is young then they are cold. $\forall x \, (\text{Young}(x) \rightarrow \text{Cold}(x))$
        \item Rule 6: If someone is young then they are nice. $\forall x \, (\text{Young}(x) \rightarrow \text{Nice}(x))$
        %\item $(\neg p \rightarrow \neg r)\wedge (\neg q \rightarrow \neg r)\wedge (\neg p \vee \neg q)\vdash \neg r $
    \end{itemize}
    
    \item \textbf{Questions:}
    \begin{itemize}
        \item Q1: Anne is cold. True/False? [Answer: T]
        \item Q2: Anne is rough. True/False? [Answer: T]  
        \item Q3: Anne is young. True/False? [Answer: T]
        \item Q4: Anne is nice. True/False? [Answer: T]
    \end{itemize}
\end{itemize}

The fact ``Anne is green or blue" combined with Rules 1 and 2 creates a classic dilemma: both possibilities lead to the same conclusion. This dilemma reasoning yields  ``Anne is cold." Rule 3 then derives ``Anne is rough" from cold, Rule 4 derives ``Anne is young" from rough, and Rule 6 derives ``Anne is nice" from young. Rule 5 creates a circular reinforcement but doesn't alter the conclusions.

The logical structure can be represented as:
\[
(G_a \lor B_a) \land (G_a \rightarrow C_a) \land (B_a \rightarrow C_a) \vdash C_a
\]

\begin{itemize}
    \item $G_a$: \text{Green}(Anne)
    \item $B_a$: \text{Blue}(Anne)
    \item $C_a$: \text{Cold}(Anne)
\end{itemize}


\subsection{Variation 1: Rule Reduction with Same Conclusions}

This variation demonstrates that removing redundant rules preserves the AI's ability to reach the same conclusions, illustrating that fewer rules can be equally effective.

\begin{itemize}
    \item \textbf{Facts:}
    \begin{itemize}
        \item Anne is green or blue
    \end{itemize}
    
    \item \textbf{Rules:}
    \begin{itemize}
        \item Rule 1: If someone is green then they are cold. $\forall x \, (\text{Green}(x) \rightarrow \text{Cold}(x))$
        \item Rule 2: If someone is blue then they are cold. $\forall x \, (\text{Blue}(x) \rightarrow \text{Cold}(x))$
        \item Rule 3: If someone is cold then they are rough. $\forall x \, (\text{Cold}(x) \rightarrow \text{Rough}(x))$
        \item Rule 4: If someone is rough then they are young. $\forall x \, (\text{Rough}(x) \rightarrow \text{Young}(x))$
        \item Rule 6: If someone is young then they are nice. $\forall x \, (\text{Young}(x) \rightarrow \text{Nice}(x))$
        % Rule 5 removed - redundant
    \end{itemize}
    
    \item \textbf{Questions:}
    \begin{itemize}
        \item Q1: Anne is cold. True/False? [Answer: T]
        \item Q2: Anne is rough. True/False? [Answer: T]
        \item Q3: Anne is young. True/False? [Answer: T]
        \item Q4: Anne is nice. True/False? [Answer: T]
    \end{itemize}
\end{itemize}

The reasoning proceeds identically to the base case: the dilemma from Rules 1-2 yields ``Anne is cold," Rule 3 yields rough, Rule 4 yields young, and Rule 6 yields nice. The removal of Rule 5 has no impact on the conclusions, demonstrating its redundancy.

\textbf{Key Insight}: AI systems that recognize this redundancy can simplify their reasoning processes without sacrificing accuracy, embodying the ``less is more" principle.

The simplified logical structure becomes:
\[
(G_a \lor B_a) \land (G_a \rightarrow C_a) \land (B_a \rightarrow C_a) \vdash C_a
\]

\[
C_a \land (C_a \rightarrow R_a) \vdash R_a
\]


\[
R_a \land (R_a \rightarrow Y_a) \vdash Y_a
\]


\[
Y_a \land (Y_a \rightarrow N_a) \vdash N_a
\]

\begin{itemize}
    \item $G_a = \mathrm{Green}(Anne)$
    \item $B_a = \mathrm{Blue}(Anne)$
    \item $C_a = \mathrm{Cold}(Anne)$
    \item $R_a = \mathrm{Rough}(Anne)$
    \item $Y_a = \mathrm{Young}(Anne)$
    \item $N_a = \mathrm{Nice}(Anne)$
\end{itemize}


\subsection{Variation 2: Rule Equivalence with Different Conclusions}

This variation replaces multiple rules with logically equivalent fewer rules, but interestingly leads to different conclusions due to the modified rule interactions.

\begin{itemize}
    \item \textbf{Facts:}
    \begin{itemize}
        %\item Anne is a person
        \item Anne is green or blue
    \end{itemize}
    
    \item \textbf{Rules:}
    \begin{itemize}
        \item Rule A: If someone is green or blue then they are cold. $\forall x \, ((\text{Green}(x) \lor \text{Blue}(x)) \rightarrow \text{Cold}(x))$
        \item Rule 3: If someone is cold then they are rough. $\forall x \, (\text{Cold}(x) \rightarrow \text{Rough}(x))$
        \item Rule 4: If someone is rough then they are young. $\forall x \, (\text{Rough}(x) \rightarrow \text{Young}(x))$
        % Rules 5 and 6 removed
    \end{itemize}
    
    \item \textbf{Questions:}
    \begin{itemize}
        \item Q1: Anne is cold. True/False? [Answer: T]
        \item Q2: Anne is rough. True/False? [Answer: T]
        \item Q3: Anne is young. True/False? [Answer: T]
        \item Q4: Anne is nice. True/False? [Answer: F]  
    \end{itemize}
\end{itemize}

Rule A directly captures the dilemma of Rules 1 and 2, leading to ``Anne is cold" with equivalent logical force. Rules 3 and 4 then proceed as before. However, without Rule 6, we cannot derive ``Anne is nice," resulting in a different conclusion for Q4.

\textbf{Key Insight}: While Rule A is logically equivalent to the combination of Rules 1 and 2, the overall rule set simplification changes the available inference paths, demonstrating that equivalence at the micro-level doesn't guarantee identical macro-level conclusions.

The logical equivalence can be shown as:
\[
(\forall x \, (G_x \rightarrow C_x) \land \forall x \, (B_x \rightarrow C_x))
\equiv 
\forall x \, ((G_x \lor B_x) \rightarrow C_x)
\]

\begin{itemize}
    \item $G_x = \mathrm{Green}(x)$
    \item $B_x = \mathrm{Blue}(x)$
    \item $C_x = \mathrm{Cold}(x)$
\end{itemize}


However, the missing Rule 6 prevents the derivation of Nice(Anne), showing that local equivalence doesn't preserve global derivability.

\subsection{Variation 3: Rule Interference with Contradictory Conclusions}

This variation adds distracting and potentially contradictory rules, testing the AI's ability to resolve conflicts and maintain coherent reasoning.

\begin{itemize}
    \item \textbf{Facts:}
    \begin{itemize}
        %\item Anne is a person
        \item Anne is green or blue
        \item Anne is not cold or not nice
    \end{itemize}
    
    \item \textbf{Rules:}
    \begin{itemize}
        \item Rule 1: If someone is green then they are cold. $\forall x \, (\text{Green}(x) \rightarrow \text{Cold}(x))$
        \item Rule 2: If someone is blue then they are cold. $\forall x \, (\text{Blue}(x) \rightarrow \text{Cold}(x))$
        \item Rule 3: If someone is cold then they are rough. $\forall x \, (\text{Cold}(x) \rightarrow \text{Rough}(x))$
        \item Rule 4: If someone is rough then they are young. $\forall x \, (\text{Rough}(x) \rightarrow \text{Young}(x))$
        \item Rule 5: If someone is young then they are cold. $\forall x \, (\text{Young}(x) \rightarrow \text{Cold}(x))$
        \item Rule 6: If someone is young then they are nice. $\forall x \, (\text{Young}(x) \rightarrow \text{Nice}(x))$
        %\item Rule 7: If someone is nice then they are not cold. $\forall x \, (\text{Nice}(x) \rightarrow \neg \text{Cold}(x))$ \quad \text{// 
        
    \end{itemize}
    
    \item \textbf{Questions:}
    \begin{itemize}
        \item Q1: Anne is cold. True/False? [Answer: F] 
        \item Q2: Anne is rough. True/False? [Answer: F] 
        \item Q3: Anne is young. True/False? [Answer: F] 
        \item Q4: Anne is nice. True/False? [Answer: F] 
    \end{itemize}
\end{itemize}


\subsection{Reasoning Process and Analysis}

    This variation introduces \textit{contradictory interference}, which challenges the AI’s ability to resolve conflicts within the logical structure. The fact ``Anne is green or blue," combined with the original chain of reasoning (Rules 1-6), leads to the conclusions that Anne is cold, rough, young, and nice. However, the fact ``Anne is not cold or not nice" introduces a conflict, as it implies that being nice would mean Anne is not cold. This creates an inconsistency that different AI systems might resolve differently:

\begin{itemize}
    \item \textbf{Conservative approach}: Detect inconsistency and withhold conclusions
    \item \textbf{Priority-based approach}: Apply rule priorities or specificity heuristics
    \item \textbf{Paraconsistent approach}: Accept some contradictions and continue reasoning
\end{itemize}

In this case, a conservative reasoning system would recognize the contradiction and potentially reject all derived conclusions, resulting in \textit{false} for all questions.

 The addition of interfering rules not only tests the AI's ability to ignore distractions but also its capacity for inconsistency detection and resolution.

The contradiction can be formally represented as:
\[
(C_a \land N_a) \land (N_a \rightarrow \neg C_a) \vdash \bot
\]

\begin{itemize}
    \item $C_a = \mathrm{Cold}(Anne)$
    \item $N_a = \mathrm{Nice}(Anne)$
    \item $\bot$ represents a value that is always false.
\end{itemize}

\subsection{Comparative Analysis}

Table~\ref{comparison_conclusion} provides a clear comparison of how different modifications to the rule set lead to distinct conclusion patterns, highlighting the sensitivity of reasoning systems to structural changes. These variations demonstrate how even minor adjustments to the logical framework can significantly impact the reasoning process and the final outcomes.

\begin{table}
\begin{tabular}{|l|c|c|c|c|}
\hline
\textbf{Variations} & \textbf{Cold} & \textbf{Rough} & \textbf{Young} & \textbf{Nice} \\
\hline
Base Example & T & T & T & T \\
Rule Reduction & T & T & T & T \\
Rule Equivalence & T & T & T & F \\
Rule Interference & F & F & F & F \\
\hline
\end{tabular}
\caption{Comparison of conclusions across variations}
\label{comparison_conclusion}
\end{table}


\end{document}
